\section{Lecture 2 (September 4th)}
\begin{defi}
(13.1) Let $K$ be a field. A field extension of $K$ is a pair $(F,i :K\rightarrow F)$ where $F$ is a field and $i$ is an injective (ring) homomorphism. We draw field extension tower diagrams like the following: 
\[
\begin{tikzcd}[row sep=2.2em]
F \arrow[d, no head] \\
K \arrow[d, no head] 
\end{tikzcd}
\]
\end{defi}
\vspace{2ex}
\begin{ex}
(13.2 $\ast$) (Algebraic extension) Let $f(t)\in {\bm Q}[t]$ be a nonzero irreducible polynomial. Then ${\bm Q}[t]\,/\,(f(t))$ is a field, and the composite map 
\[{\bm Q}\rightarrow {\bm Q}[t]\rightarrow {\bm Q}[t]\,/\,(f(t))\]
is a field extension. This is the general way we would be obtaining field extensions.
\[
\begin{tikzcd}[row sep=2.2em]
{\bm Q[t]}\,/\,(f(t)) \arrow[d, no head] \\
{\bm Q} \arrow[d, no head] 
\end{tikzcd}
\]
\end{ex}
\vspace{2ex}
\begin{proof}
The set ${\bm Q}[t]\,/\,(f(t))$ is a maximal ideal and since $f(t)$ is a nonzero irreducible element of ${\bm Q}[t]$ which is principle ideal domain. Therefore, the quotient ${\bm Q}[t]\,/\,(f(t))$ is a field. 
\end{proof}
\vspace{2ex}
\begin{ex}
(13.3) (Transcendental extension) Let ${\bm Q}(t)$ denote the field of fractions of ${\bm Q}[t]$. Then the composition map ${\bm Q}\rightarrow {\bm Q}[t]\rightarrow {\bm Q}(t)$ determines a field extension.
\end{ex}
\vspace{2ex}
\begin{defi}
(13.6) The characteristic of a field $K$ ($\mathop{\mathrm{char}}(K)$) is the smallest nonnegative generator of the kernel of the unique ring homomorphism $\phi :{\bm Z}\rightarrow K$. 
\end{defi}
\vspace{2ex}
\begin{lem}
(13.7) Let $F\,/\,K$ be a field extension. Then 
\[\mathop{\mathrm{char}}(K)=\mathop{\mathrm{char}}(F)\]
\end{lem}
\vspace{2ex}
\begin{lem}
(13.8) Let $F$ be a field. Then $F$ contains a unqiue copy of ${\bm Q}$ or ${\bm Z}\,/\,p{\bm Z}$, where $p$ is prime.
\end{lem}
\vspace{2ex}
\section{13.2 Finite Extensions and the Degree of an Extension}
\begin{defi}
(13.10) The degree of a field extension $F/K$, denoted $[F:K]$, is the dimension of $F$ as a vector space over $K$.
\[[F:K]=\mathrm{dim}_{K}(F)\]
\end{defi}
\vspace{2ex}
\begin{rmk}
If $i:K\rightarrow F$ is a field extension, then we can view $F$ as a $K$-vector space through the map $i$.
\end{rmk}
\vspace{2ex}
\begin{defi}
(13.11) A field extension is finite if $[F:K]<\infty $.
\end{defi}
\vspace{2ex}
\begin{defi}
(algebraic and transcendental elements) Let $F/K$ be a field extension, and let $\alpha \in F$. We will say that $\alpha $ is algebraic over $K$ if there exists a nonzero polynomial $f(t)\in K[t]$ such that $f(\alpha )=0$. The element $\alpha \in F$ is transcendental if not algebraic.
\end{defi}
\vspace{2ex}
\begin{ex}
Consider
\[
\begin{tikzcd}[row sep=2.2em]
{\bm C} \arrow[d, no head] \\
{\bm R} \arrow[d, no head] 
\end{tikzcd}
\]
\begin{itemize}
\item[(i)] Is $i\in {\bm C}$ algebraic over ${\bm R}$? Yes! $i$ is a zero of $t^2+1\in {\bm R}[t]$.
\item[(ii)] $\sqrt{2}\in {\bm R}$ is algebraic over ${\bm Q}$
\item[(iii)]  Also, $\pi, e \in {\bm C}$ is transcendental over ${\bm Q}$. 	
\item[(iv)] Every element $\alpha $ of $K$ is algebraic over $K$ as we have the polynomial 
\[t-\alpha \in K[t]\]
\end{itemize}
\end{ex}
\vspace{2ex}
\begin{defi}
(13.20) Let $F/K$ be a field extension. If $\alpha \in F$ is algebraic, the minimal polynomial of $\alpha $ over $K$ is the unique monic irreducible polynomial $p(t)\in K[t]$ such that $f(\alpha )=0$.
\end{defi}
\vspace{2ex}
\begin{proof}
We show that the minimal polynomial exists and is unique. We will achieve this by showing that it is exactly the generator of the kernel of the evaluation map. First consider the evaluation homomorphism
\[\mathrm{ev}_{\alpha }:K[t]\rightarrow F\]
where
\[f(t)\mapsto f(\alpha )\qquad\&\qquad \sum ^{n}_{i=0}a_{i}t^{i}\mapsto \sum ^{n}_{i=0}a_{i}\alpha ^{i}\]
Observe that it is indeed a homomorphism as it preserves addition, multiplication, and the identity element. Then, notice the following:

\bigskip

(I) The kernal of the map is automatically an ideal, and is of the form $(p(t))$ as $K[t]$ is a principle ideal domain.

\bigskip

(II) By the 1st isomorphism theorem, the following isomorhpism holds
\[K[t]\,/\,\mathrm{ker}(\mathrm{ev}_{\alpha })=K[t]\,/\,(p(t))\;\cong\; \mathrm{im}(\mathrm{ev}_{\alpha })=K[\alpha ]\subset F\]
This tells us that $K[\alpha ]$ is a subring of $F$, $K[\alpha ]$ is therefore a subring and an integral domain, and ultimately that $(p(t))$ is a prime ideal and that $p(t)$ is an irreducible element (we know that it is nonzero because $\alpha $ is algebraic).

\bigskip

(III) As generators are unique up to multiplication by a unit, we might as well suppose that the polynomial $p(t)$ is monic.

\bigskip

(IV) We lastly need to prove uniqueness. Suppose that there's another irreducible monic polynomial such that $q(t)=0$. As $q(t)\in \mathrm{ker}(\mathrm{ev}_{\alpha })=(p(t))$, $q(t)=p(t)\cdot r(t)$ and as $q(t)$ is irreducible, $r(t)$ has to be a unit. Lastly, $r(t)$ has to be 1 as both $q(t)$ and $p(t)$ are monic polynomials. 
\[r(t)\in (K[t])^{\times }=K^{\times }=K\qquad\&\qquad p(t)=q(t)\]

\end{proof}
\vspace{2ex}

